Abstract / Task Definition

1. INTRODUCTION
    1.1 Background and Motivation
        1.1.1 Automotive Software Evolution and Complexity
        1.1.2 Security Challenges in Modern Systems
        1.1.3 CI/CD/CT Pipeline Context
    1.2 Problem Statement
    1.3 Research Questions and Objectives
        1.3.1 Primary Research Questions
        1.3.2 Secondary Research Questions
    1.4 Thesis Organization

2. LITERATURE REVIEW
    2.1 Traditional AI Approaches in Software Testing
    2.2 Large Language Models: Evolution and Code Generation
    2.3 Foundations to AI-Guided Fuzzing Techniques
    2.4 AI-Enhanced Fuzzing: State-of-the-Art
        2.4.1 LLM-Based Test Case Generation
        2.4.2 Neural Network-Guided Fuzzing
    2.5 CI/CD Pipeline Security Integration
    2.6 Automotive Software Development and Safety Standards
    2.7 Research Gaps and Thesis Positioning

3. METHODOLOGY AND FRAMEWORK
    3.1 Conceptual Framework: AI-Driven White-Box Fuzzing
    3.2 Technical Architecture Design
    3.3 Research Strategy: Three-Phase Approach
        3.3.1 Phase 1: Local LLM Evaluation
        3.3.2 Phase 2: Model Optimization with LoRA
        3.3.3 Phase 3: Enterprise CI/CD Integration
    3.4 Evaluation Methodology
        3.4.1 Metrics and Performance Indicators
        3.4.2 Validation and Verification Approach

4. IMPLEMENTATION
    4.1 Toolchain Selection and Rationale
    4.2 Phase 1: Local LLM Evaluation Setup
        4.2.1 Benchmarked Models
        4.2.2 Target Repositories
        4.2.3 Evaluation Environment
    4.3 Phase 2: Model Optimization with LoRA Fine-Tuning
    4.4 Phase 3: Enterprise CI/CD Integration
        4.4.1 Architectural Challenges
        4.4.2 Self-Hosted Runner Solution

5. EXPERIMENTAL RESULTS
    5.1 Experimental Setup
        5.1.1 Target Selection and Criteria
        5.1.2 Hardware and Software Configuration
    5.2 LLM Fuzz Driver Generation Results
        5.2.1 Successful Models: Performance Data
        5.2.2 Unsuccessful Models: Critical Findings
        5.2.3 Performance Across Repositories
    5.3 Model Optimization Results
        5.3.1 LoRA Fine-Tuning Efficiency
        5.3.2 Comparative Analysis
    5.4 Economic Analysis and Resource Metrics

6. DISCUSSION AND CONCLUSION
    6.1 Interpretation of Results
        6.1.1 What the Data Reveals
        6.1.2 Unexpected Findings: Network Architecture
    6.2 Addressing Research Questions
        6.2.1 Primary Research Question Analysis
        6.2.2 Secondary Research Questions Analysis
    6.3 Limitations and Constraints
        6.3.1 Technical Limitations
        6.3.2 Methodological Boundaries
    6.4 Implications for Practice
        6.4.1 Automotive Industry Impact
        6.4.2 Enterprise CI/CD Challenges
    6.5 Future Research Directions
    6.6 Summary of Findings and Contributions
    6.7 Conclusion

Bibliography
Declaration / Erklärung
Appendices (Optional)